\section{Method}
\label{sec:method}

We change one link of \Cref{fig:pipeline} at a time.
This section specifies the four interventions---representation,
injection, trajectory measurement, and sampling-time
guidance---before any SR table.
Training hyperparameters that never change are in
\Cref{tab:constants}; checkpoint pairing is in \Cref{tab:ckpts}.

\subsection{Stage A: reconstruction VAE (\vaeone)}
\label{sec:method-vae1}

\vaeone{} is a convolutional $\beta$-VAE~\citep{kingma2014auto}
trained only to reconstruct HR CelebA faces.
It has no SR term.
The encoder outputs $(\mu,\log\sigma^2)$; training uses
$z=\mu+\sigma\odot\eps$ with $\eps\sim\Ncal(0,I)$.
The loss is mean-reduced MSE plus a mean-reduced KL, so the KL
scale stays stable at latent size $4\times 32\times 32$:
\begin{equation}
  \label{eq:vae1}
  \begin{aligned}
    \Lcal_{\mathrm{VAE\text{-}1}}
    \;=\;
    &\underbrace{\lVert \Dec(z)-\xhr\rVert_2^2}_{\Lcal_{\mathrm{recon}}} \\
    &+\beta\,
    \underbrace{\Ebb\bigl[-\tfrac12\bigl(1+\log\sigma^2-\mu^2-\sigma^2\bigr)\bigr]
    }_{\mathrm{KL}},
  \end{aligned}
\end{equation}
with $\beta=10^{-4}$ and
$\mathrm{KL}=\mathrm{KL}\bigl(q(z\mid\xhr)\,\Vert\,\Ncal(0,I)\bigr)$.
We train for $50$ epochs (Adam $10^{-4}$, batch $64$) and freeze
epoch~$50$ for every later experiment.
Eval reconstructions use $\mu$, not a posterior sample.

\subsection{Stage B: SR-aware alignment (\vaesr)}
\label{sec:method-vaesr}

\vaesr{} is \vaeone{} fine-tuned so the LR encoder mean moves toward
the HR encoder mean without rebuilding the autoencoder
(\Cref{fig:stage-vae}).
Initialize from \vaeone{} epoch~$50$.
On each batch encode HR and the bicubic-upsampled LR with the
\emph{same} encoder
\begin{equation}
  \label{eq:mu-pair}
  \mu_{\mathrm{hr}},\log\sigma^2_{\mathrm{hr}} \;=\; \Enc(\xhr),\qquad
  \mu_{\mathrm{lr}} \;=\; \Enc(\xup),
\end{equation}
and optimize
\begin{equation}
  \label{eq:vaesr}
  \begin{aligned}
    \Lcal_{\mathrm{VAE\text{-}SR}}
    \;=\;
    &\Lcal_{\mathrm{recon}}(\xhr,\Dec(z_{\mathrm{hr}})) \\
    &+\beta\,\mathrm{KL}\bigl(q(z\mid\xhr)\,\Vert\,\Ncal(0,I)\bigr) \\
    &+\lamal \lVert \mu_{\mathrm{lr}} - \sg(\mu_{\mathrm{hr}}) \rVert_2^2.
  \end{aligned}
\end{equation}
Stop-gradient $\sg(\cdot)$ on $\mu_{\mathrm{hr}}$ is required: without
it the HR code can collapse toward the blurry LR code; with it, only
the LR mean is pulled.
$\beta$ is unchanged ($10^{-4}$).
$\lamal=10^{-3}$ because unweighted align MSE starts near $0.58$;
$10^{-3}$ keeps the term on the order of recon MSE.
We fine-tune for $20$ epochs (Adam $10^{-4}$, batch $32$) and freeze
the result.
A $\lamal$ sweep is \emph{not} part of the method: cosine$(\zlr,\zhr)$
plateaus after epoch~$1$, so the next causal question is consumption,
not more alignment.

\begin{figure}[htbp]
  \centering
  \begin{tikzpicture}[
      font=\scriptsize,
      blk/.style={draw, rounded corners=1.5pt, align=center, inner sep=3pt,
                  minimum height=8.5mm, anchor=center},
      arr/.style={-{Latex}, thick, shorten >=1pt, shorten <=1pt}
    ]
    \matrix[
      ampersand replacement=\&,
      column sep=6.5mm,
      row sep=1.45cm,
      cells={anchor=center, inner sep=0pt}
    ] {
      \node[blk, fill=gray!12] (hr) {$\xhr$}; \&
      \node[blk, fill=blue!12] (ehr) {$\Enc$}; \&
      \node[blk, fill=blue!12] (muh) {$\mu_{\mathrm{hr}}$}; \&
      \node[blk, fill=orange!15] (dec) {$\Dec(z_{\mathrm{hr}})$}; \\
      \node[blk, fill=gray!12] (lr) {$\xup$}; \&
      \node[blk, fill=blue!12] (elr) {$\Enc$}; \&
      \node[blk, fill=blue!12] (mul) {$\mu_{\mathrm{lr}}$}; \&
      \node[blk, fill=green!12] (aln) {$\lamal\lVert\mu_{\mathrm{lr}}$\\
        ${}-\sg(\mu_{\mathrm{hr}})\rVert_2^2$}; \\
    };
    \draw[arr] (hr) -- (ehr);
    \draw[arr] (lr) -- (elr);
    \draw[arr] (ehr) -- (muh);
    \draw[arr] (elr) -- (mul);
    \draw[arr] (muh) -- (dec);
    \draw[arr] (mul) -- (aln);
    % Stop-grad: label sits on the arrow, not floating above it.
    \draw[arr] (muh.south east) to[out=-75, in=90]
      node[pos=0.5, fill=white, inner sep=1.2pt, font=\tiny] {$\sg$} (aln.north);
    \node[font=\tiny\itshape, text=black!55, above=0.08cm of hr] {\vaeone{} path (kept)};
    \node[font=\tiny\itshape, text=black!55, below=0.08cm of lr] {new SR term};
  \end{tikzpicture}
  \caption{\vaesr{} fine-tune.
  HR recon+KL is the \vaeone{} loss~\eqref{eq:vae1}.
  Alignment pulls $\mu_{\mathrm{lr}}$ toward a stop-grad $\mu_{\mathrm{hr}}$.
  Soft-decode later scores $\Dec(\mu_{\mathrm{lr}})$, not this training grid.}
  \label{fig:stage-vae}
\end{figure}

\subsection{Stage C: concat latent DDPM}
\label{sec:method-concat}

Given a frozen VAE, we train a noise-prediction UNet on
$z_{\mathrm{hr}}$ with channel-concat conditioning~\eqref{eq:concat}:
eight input channels, four output channels, same skeleton and
$50$-epoch recipe for both VAEs.
The training loss is standard DDPM noise MSE~\citep{ho2020denoising}.
Only the frozen VAE changes between the \vaeone{} concat run
(Phase~8) and the \vaesr{} concat run (Q2).
We do not resume a concat checkpoint into FiLM (channel layout
differs).

\subsection{Stage D: spatial FiLM}
\label{sec:method-film}

To test whether concat is a weak reader of an already-better
$\zlr$, we replace~\eqref{eq:concat} with spatial FiLM
\citep{perez2018film} at every UNet scale.
The noisy latent stays four channels.
At a feature map $h$ of matching (or bilinearly resized) resolution,
\begin{equation}
  \label{eq:film}
  h' \;=\; \bigl(1+\gamma(\zlr)\bigr)\odot h \;+\; \beta(\zlr),
\end{equation}
where $\gamma,\beta$ come from a $1{\times}1$ convolution that is
\emph{zero-initialized}, so the block is identity at step~$0$.
Cross-attention is not used (\S\ref{sec:related}).

\paragraph{One-sided $2{\times}2$.}
A full interaction would train FiLM on both VAEs.
With one new DDPM, we train FiLM+\vaesr{} and interpret it as a
one-sided test, written down \emph{before} seeing numbers:
if SR still matches concat+\vaesr, skip FiLM+\vaeone{} and treat
concat-as-reader as not the bottleneck;
if FiLM+\vaesr{} clearly beats concat+\vaesr, then train FiLM+\vaeone{}
before claiming an interaction.
Same $50$ epochs, same paired eval as Q2.

\subsection{Stage E: reverse-chain diagnostics}
\label{sec:method-traj}

To localize where a better $\zlr$ is or is not used, we record, at
every scheduler index $t$,
$\cos(\zhat(t),\zlr)$ and $\lVert\zhat(t)-\zlr\rVert$ under both
frozen VAEs and their matched concat DDPMs, plus the image-space
scores of $\Dec(\zhat(t))$ versus $\xhr$.
Soft-decode $\Dec(\zlr)$ is not $\Dec(\zhat(t))$ at a copy peak;
those are reported separately.

\subsection{Stage F: DPS-style LR-consistency guidance}
\label{sec:method-guide}

Guidance is inference only: frozen \vaesr{} + Q2 concat, no extra
training.
The convention is DPS-like~\citep{chung2023diffusion} and is fixed
before any guided SR number: gradient with respect to $\zt$,
correction applied to the already-sampled $z_{t-1}$.

At each reverse index $t$:
\begin{enumerate}
  \item $\epshat=\mathrm{UNet}(\zt,\zlr,t)$, \emph{detached}
        (no UNet backward).
  \item $\zhat$ as in~\eqref{eq:z0hat}.
  \item $z_{t-1}\sim p(z_{t-1}\mid \zt,\epshat)$ (unguided DDPM step).
  \item If $t$ is in the window,
\end{enumerate}
\begin{equation}
  \label{eq:guide}
  \begin{aligned}
    \Lcal_g
    &=
    \mathrm{mean}_{c,h,w}\bigl\lVert \Dec(\zhat)-\Dec(\zlr)\bigr\rVert_2^2, \\
    g
    &=
    \nabla_{\zt}\Lcal_g, \\
    z_{t-1}
    &\leftarrow
    z_{t-1}-\lamg\, g.
  \end{aligned}
\end{equation}
$\Dec(\zlr)$ is cached once per image.
Decoder weights stay frozen; the backward runs through the
$\zhat$ affine and $\Dec$ only.
Windows (scheduler index; $t{=}0$ is clean):
baseline $=$ none;
early $=$ $t\in[0,800]$ every step;
late $=$ $t\in[0,500]$ every step.
A budget skip \texttt{guide\_every}$>1$ is implemented but not used
in the reported runs.

\paragraph{Choosing $\lamg$ from $R(t)$, not from SR PSNR.}
Define the relative step size
\begin{equation}
  \label{eq:Rt}
  R(t) \;=\; \frac{\lVert g\rVert_2}{\lVert z_{t-1}-z_t\rVert_2}.
\end{equation}
A sanity sweep at $\lamg\in\{10^{-3},10^{-2},10^{-1},1\}$ gives
$R(t)<0.01$ everywhere: $\lamg{=}1$ is a no-op.
$R$ is linear in $\lamg$.
\Cref{tab:Rt} scales the $\lamg{=}1$ column.
$\lVert g\rVert/\lVert\Delta z\rVert$ drops by about $60\times$ from
$t{=}800$ to $t{=}500$, so \emph{one scalar cannot put every
checkpoint in a $0.1$--$0.5$ band}.
We take $\lamg{=}200$ as a compromise: in-band at $t{=}650$
($R{\approx}0.18$); above $0.5$ but below a $2$--$3$ dominate
ceiling at $t{=}800$ ($R{\approx}1.16$); small but non-zero in the
late window ($R(500){\approx}0.04$).
$\lamg{=}400$ would make the late window fairer but
$R(800){\approx}2.3$.
A late PSNR-null is not treated as localization failure unless
$R(t)$ for $t\le 500$ is clearly $>0$.

\begin{table}[htbp]
  \centering
  \caption{Relative guidance size $R(t)$~\eqref{eq:Rt}, scaled from
  a $\lamg{=}1$ measurement. Bold: operating column.}
  \label{tab:Rt}
  \small
  \begin{tabular}{@{}rcccc@{}}
    \toprule
    $t$ & $R(1)$ & $R(100)$ & $\mathbf{R(200)}$ & $R(400)$ \\
    \midrule
    $800$ & $0.0058$ & $0.58$ & $\mathbf{1.16}$ & $2.32$ \\
    $650$ & $0.0009$ & $0.09$ & $\mathbf{0.18}$ & $0.36$ \\
    $500$ & $0.0002$ & $0.02$ & $\mathbf{0.04}$ & $0.08$ \\
    $300$ & $0.0001$ & $0.01$ & $\mathbf{0.02}$ & $0.04$ \\
    $100$ & $0.0001$ & $0.01$ & $\mathbf{0.02}$ & $0.04$ \\
    \bottomrule
  \end{tabular}
\end{table}

\paragraph{Pre-registered calls (before guided PSNR).}
PSNR$\uparrow$ with LPIPS $\approx 0.068$ is a strong positive;
PSNR$\uparrow$ with LPIPS toward $0.119$ is over-guidance.
If both windows work, Stage~2 is a late-window dose
$\lamg\in\{50,200,800\}$ only---no combined schedule.
The $n{=}256$ confirmation (val\_index $0..255$, same checkpoints,
late window) is declared before looking at those numbers, and
survives iff (a)~$\ge 90\%$ of images have $\Delta\mathrm{PSNR}>0$ at
each $\lamg$, (b)~mean $\Delta$LPIPS is monotone
$50<200<800$ with $\lamg{=}800$ worse than unguided and
$\lamg{=}50$ not worse, (c)~PSNR CIs exclude $0$ at all three
$\lamg$.
The first line of each job must reprint
$\mathrm{PSNR}(\Dec(\zlr),\xhr)\approx 28.48$ on indices $0..63$;
otherwise the wrong VAE is loaded and the run is discarded.
